\documentclass[11pt]{article}
\usepackage[T1]{fontenc}
\usepackage[utf8]{inputenc}
\usepackage{lmodern,amsmath,amssymb,amsthm}
\usepackage[margin=1in]{geometry}
\usepackage[hidelinks]{hyperref}
\newtheorem{theorem}{Theorem}
\newtheorem{lemma}[theorem]{Lemma}
\newtheorem{corollary}[theorem]{Corollary}
\newcommand{\N}{\mathbb N}\newcommand{\Z}{\mathbb Z}
\newcommand{\R}{\mathbb R}\newcommand{\Q}{\mathbb Q}
\newcommand{\C}{\mathbb C}\newcommand{\g}{\mathfrak g}
\newcommand{\Sym}{\operatorname{Sym}}\newcommand{\poly}{\operatorname{poly}}
\title{Scaling symbols and symmetric multilinearisation with the factorial retained}
\author{Author proposal: /root/astra\_v14\_baseline\_compiler}
\date{7 November 2026}
\begin{document}\maketitle
\noindent\textbf{Status and scope.} This is a new proof proposal, not an
independent acceptance. It supplies GTZ-LP05 and the scaling assertion
E.8(v), including the divisibility used here. Its project input is the exact
tautification theorem in \cite{TAUT}, still pending independent review.
It does not prove the remaining correction units or the inverse theorem.
The unpublished erratum motivates the factorial correction but supplies no
black-box inference. All source and project pins are in the accompanying
contracts. Protocol v1.4's available/needed/connecting-argument distinctions
are recorded there.

\section{Conventions and preliminary constructions}\label{sec:setup}
We use the rational, connected, simply connected nilpotent groups,
polynomial maps, finite-dimensional unit-sphere-valued nilcharacters and
symbol equivalence of \cite[Sections 3--6 and Appendix E]{GTZ}.
For a vector-valued function $u$, $u^{\otimes r}$ means its ordinary tensor
power if $r\geq0$, with $u^{\otimes0}=1$, and
$\overline u^{\otimes(-r)}$ if $r<0$.
Conjugation is always written explicitly. Equivalence at multidegree $d$
means that every component of $u\otimes\overline v$ is a nilsequence with
support in $\{a:a\leq d,\ a\ne d\}$; all inequalities on multiindices
are coordinatewise. In the degree case this is degree at most $d-1$.
We use the group law on symbols proved in published E.7 and E.8(i)--(iii).
Their proofs use tensor products, conjugation, contraction of a unit vector
with its conjugate, and a rational top-group quotient. With tautification,
the top group is central, so those proofs apply to the present objects.

Initially all data are ordinary, and all integers $d,k,q$ are fixed standard
integers. The last section states exactly how the constructions pass to
limits and parameter families. A multidegree nilcharacter on a limit subset
comes with a representing global polynomial orbit and test function; we
use that representation to extend it, and only then restrict back.

Apply \cite[Theorem 1 and Corollary 2]{TAUT} to the given representation.
We can work on a taut group and with an exactly equal represented sequence
whose polynomial map $g$ satisfies $g(0)=1$. Positive groups and the
vertical character are unchanged. For a degree filtration $G_0=G_1$;
for multidegree $G_0$ is generated by the $G_{e_i}$. In particular $G_d$
is central when the support is contained below one positive $d$.
Write $\g_a=\log G_a$ with its rational structure induced by $\Gamma$.

\begin{lemma}[Monomial coordinates]\label{lem:monomial}
For multidegree support $a\leq d$, the normalised map has an expression
\[
 g(n)=\prod_{0<a\leq d}\exp(X_a n^a),\qquad X_a\in\g_a,
 \quad n^a=\prod_i n_i^{a_i}.                 \tag{1}
\]
Fix any order refining total weight. For a degree-$D$ map on $\Z^k$,
replace $a\leq d$ by $0<|a|\leq D$ and require $X_a\in\g_{|a|}$.
Every such ordered product is polynomial for the indicated filtration.
\end{lemma}
\begin{proof}
For a degree-$D$ input, use the multidegree target $\g_a=\g_{|a|}$:
coordinate derivatives are particular degree derivatives, so the input
is polynomial for this target. Published Lemma B.9 gives ordered
binomial monomials with coefficient
at index $a$ in $G_a$. Expand the finite BCH formula. A product of scalar
binomial polynomials with labels $a_1,\ldots,a_r$ has monomials $n^b$
only for $b\leq a_1+\cdots+a_r$. Its Lie-bracket coefficient belongs to
$\g_{a_1+\cdots+a_r}\subseteq\g_b$. Thus
$\log g(n)=\sum_{0<b\leq d}Y_b n^b$, $Y_b\in\g_b$; terms outside the
support vanish and the constant is zero. To obtain (1), process increasing
total monomial weight: choose $X_b$ to match the coefficient at $n^b$ in
the BCH logarithm of the partial ordered product. All commutators of
previous factors have larger total monomial weight, and their coefficients
at $b$ belong to $\g_b$. Consequently the adjustment is in $\g_b$ and
does not change any previously matched weight. The finite support makes
this algorithm terminate. The argument with $\g_{|b|}$ is identical.
For the converse each factor is a polynomial map: in its one-parameter
abelian subgroup an iterated difference of $n^a$ vanishes if its coordinate
order exceeds $a$, and otherwise the coefficient stays in
$\g_a\subseteq\g_b$ at derivative index $b$. Apply the product theorem
\cite[Corollary B.4]{GTZ}. This use does not invoke E.8, E.9 or GI.
\end{proof}

We will use rational connected subgroups and rational central quotients.
Their generated subgroups and central quotients are the geometric
construction proved in \cite[Section 2]{LP13}. We spell out the additional
rational-image check used below. For each rational $G_a$, choose finitely
many elements $\gamma_{a,j}$ of $\Gamma\cap G_a$ whose logarithms span
$\g_a$. Such a choice exists: the real span of the logarithms of a
nilpotent subgroup is a Lie algebra. Indeed conjugation by integer powers
of any of its elements gives a polynomial in that integer, and finite
differences extract its iterated brackets; BCH then gives the assertion.
Its connected subgroup contains the lattice, so cocompactness forces it
to be $G_a$, by the homogeneous-space argument of \cite[Section 2]{LP13}.
The same polynomial/BCH argument gives the rational Lie structure
$\operatorname{span}_{\Q}\log\Gamma$; all filtration spaces are rational
for this structure.

For the scaling image below, the vectors
$(q^{|a|}\log\gamma_{a,j},\log\gamma_{a,j})$ exponentiate to the actual
lattice elements $(\gamma_{a,j}^{q^{|a|}},\gamma_{a,j})$. For the
multilinear image, multiply each $\log\gamma_{a,j}$ by an integer $N$
so that its rational second component in (8) lies in the lattice $L$
constructed next. Its first component then exponentiates to
$\gamma_{a,j}^N$. These finitely many lattice elements have logarithms
spanning each image Lie algebra $\mathfrak k_b$. The finite-generator
fact in \cite[p.~1243]{GTZ} makes their generated discrete subgroup
cocompact in $K_b$. Thus $K_b$ is rational for the actual product lattice,
not merely for an abstract rational form. This also checks $b=0$.

The following explicit lattice choice records
the extra finite-symmetry requirement needed below. In a rational
nilpotent Lie algebra with a finite group $P$ of rational automorphisms,
choose a full rational lattice $L_0$ and replace it by
$\sum_{\sigma\in P}\sigma L_0$. Clear rational bracket denominators and
the finitely many BCH denominators by replacing $L_0$ by $L=M L_0$.
Then BCH$(L,L)\subseteq L$ and $-L=L$. The group $\exp L$ is discrete.
It is cocompact: an adapted rational central-series basis, after integer
scaling, lies in $L$; successive rounding of its coordinates gives a
compact set of coset representatives. The same construction makes all
specified rational filtration subspaces rational. If a rational linear
functional on a specified rational central subspace is prescribed,
increase $M$ to make it integral on $L$ intersected with that subspace.
This gives a $P$-invariant lattice, not just a rational structure invariant
under $P$. Each choice uses only fixed geometric data.

\section{Scaling and divisibility on the actual orbit}\label{sec:scale}
\begin{theorem}[Corrected E.8(v)]\label{thm:scale}
Let $d\ne0$, and let $\chi(n)=F(g(n)\Gamma)$ have multidegree at most $d$.
For every standard integer $q$,
\[
 [\chi(qn)]=q^{|d|}[\chi(n)] .                \tag{2}
\]
For degree $D>0$ on $\Z^k$, the exponent is $D$. The equivalence is on
the whole domain, hence on every indicated restriction. Its geometric
and Lipschitz bounds are independent of the polynomial coefficients.
\end{theorem}
\begin{proof}
Use (1). Form the finite positive Rees Lie algebra
\[
 \mathcal R=\bigoplus_{0<a\leq d}\g_a t^a,
 \qquad[Xt^a,Yt^b]=[X,Y]t^{a+b},
\]
where a component outside the downset is zero. This is a Lie algebra:
it is a subalgebra of $\g\otimes\R[t_1,\ldots,t_k]$ modulo the ideal
spanned by monomials outside the downset. Define the rational Lie map
\[
 A_q(Xt^a)=(q^{|a|}X,X)\in\g\oplus\g.
\]
Its bracket identity is exact because
$q^{|a|}q^{|b|}=q^{|a+b|}$, including negative and zero $q$.
Let $\mathfrak k_b$ be the image of
$\bigoplus_{a\geq b,\ a\ne0}\g_a t^a$, and put $K_b=\exp\mathfrak k_b$.
These are rational, descending, and
$[\mathfrak k_b,\mathfrak k_c]\subseteq\mathfrak k_{b+c}$.
Also $\mathfrak k_0=\sum_i\mathfrak k_{e_i}$, so the filtration is taut.
The ambient group is $K_0$, with lattice $K_0\cap\Gamma^2$.
Every factor of
\[
 (g(qn),g(n))=\prod_a\exp\bigl(n^a A_q(X_a t^a)\bigr)
\]
is polynomial at index $a$, and hence so is their product. The top group
has Lie algebra
$\mathfrak k_d=\{(q^{|d|}X,X):X\in\g_d\}$.
For $c=q^{|d|}$, the restricted function
\[
 (x,y)\longmapsto F(x)\otimes\overline{F(y)^{\otimes c}}
                                                        \tag{3}
\]
is invariant under $K_d$: the two vertical phases are $c\eta(X)$ and
$-c\eta(X)$. This calculation uses $F^{\otimes c}$, with the signed
convention above, on the same second coordinate. The rational central
quotient $K_0/K_d$ has the quotient filtration with its top group zero.
Thus (3) along the displayed orbit is a lower-degree nilsequence, which
is precisely (2). Restriction and descent preserve a fixed Lipschitz
bound since the geometric maps and compact quotient are fixed.
For degree $D$ use
$\bigoplus_{0<|a|\leq D}\g_{|a|}t^a$ and the filtration by $|a|\geq j$.
The top image is the same graph with exponent $D$; all other steps agree.
At $q=0$ the first orbit coordinate is $1$, and (3) is the constant
$F(1)$; no $0^0$ or division is used.
\end{proof}

\begin{corollary}[Divisibility, with the representative specified]\label{cor:divide}
For positive degree and every standard $Q\geq1$, there is a nilcharacter
$\psi$ on the same domain with $[\chi]=Q[\psi]$.
\end{corollary}
\begin{proof}
In (1) replace each $X_a$ by $Q^{-|a|}X_a$ and call the resulting map
$g_Q$. It is polynomial on the same group and lattice. With the same
$F$ put $\xi(n)=F(g_Q(n)\Gamma)$. Then $g_Q(Qn)=g(n)$ exactly, factor
by factor in the same order. Theorem~\ref{thm:scale} gives
$[\chi]=Q^D[\xi]$, where $D=|d|$ or the degree bound, and
$\psi=\xi^{\otimes Q^{D-1}}$ gives the asserted equation.
The test function still has unit norm; its frequency and finite target
dimension change only by the displayed standard tensor power. In
particular this is divisibility of the actual symbol, not division of
an arbitrarily selected multilinear representative after construction.
\end{proof}

\section{Raw multilinearisation and the factorial graph}\label{sec:raw}
Let $d=(d_1,\ldots,d_k)\ne0$, $D=|d|$, and divide $\{1,\ldots,D\}$
into ordered blocks $B_i$ of size $d_i$. For a subset $J$ put
$w(J)=(|J\cap B_i|)_i$ and $a!=\prod_i a_i!$; set $C=d!$.
Construct
\[
 \widetilde\g=\bigoplus_{\varnothing\ne J\subseteq[D]}\g_{w(J)}\epsilon_J,
 \qquad
 [X\epsilon_J,Y\epsilon_L]=
 \begin{cases}[X,Y]\epsilon_{J\cup L},&J\cap L=\varnothing,\\
 0,&J\cap L\ne\varnothing.
 \end{cases}                                             \tag{4}
\]
This is a subalgebra of $\g\otimes
\R[\epsilon_1,\ldots,\epsilon_D]/(\epsilon_1^2,\ldots,\epsilon_D^2)$;
Jacobi therefore holds. There is no empty component: normalisation
$g(0)=1$ is what permits its removal. Put
\[
 \widetilde\g_b=
 \bigoplus_{1_J\geq b}\g_{w(J)}\epsilon_J.                 \tag{5}
\]
These are descending rational Lie subalgebras, their bracket has the
required index sum, and they vanish when any $b_i>1$. The zero group
is generated by the first-coordinate groups since each nonempty $J$
contains a coordinate. Thus (5) is taut. Its top group is central.
The block permutation group $P=\prod_i S_{d_i}$ preserves (4), permutes
the groups in (5) by the same permutation of their indices, and
acts trivially on the top copy $\g_d\epsilon_{[D]}$.
Choose a $P$-invariant lattice $\widetilde\Gamma$ as in Section
\ref{sec:setup}, making
\[
 \widetilde\eta(\exp(X\epsilon_{[D]}))=\eta(\exp X)
                                                        \tag{6}
\]
integral on its top lattice. This is possible since $\eta$ is integral
on $\Gamma\cap G_d$ and hence rational on the rational top Lie algebra.

Choose a Lipschitz unit-sphere-valued function $\widetilde F$ with frequency
$\widetilde\eta$. Here is the needed existence construction, with no
invariance under $P$ assumed. The central torus acts freely on
$\widetilde G/\widetilde\Gamma$: its stabilizer is precisely the top
lattice. On finitely many local trivialisations of this torus bundle,
take the character (6) times a smooth base bump. Choose the bumps so
that their supports cover the compact base and they vanish near the
chart boundaries. The vector of these functions has no common zero;
divide by its norm. This gives a globally Lipschitz unit vector of the
required frequency, with a fixed finite dimension. It is the
partition-of-unity construction underlying published (6.3).

For the coefficients of (1), define
\[
 \widetilde g(z)=\prod_{0<a\leq d}
 \exp\left(a!\sum_{w(J)=a}z_J X_a\epsilon_J\right),
 \qquad z_J=\prod_{j\in J}z_j,\qquad
 \widetilde\chi(z)=\widetilde F(\widetilde g(z)\widetilde\Gamma).
                                                        \tag{7}
\]
Each factor is polynomial for (5). To see this without commuting distinct
coefficients, fix $a$; the summands in its exponent commute because their
Lie coefficients are the same $X_a$, so (7)'s factor is a product of the
individual monomial factors. Then use product closure. Thus
$\widetilde\chi$ has multidegree at most $(1,\ldots,1)$.
Let $\Delta n$ repeat $n_i$ throughout $B_i$.

\begin{lemma}[Raw diagonal identity]\label{lem:raw}
The diagonal $\widetilde\chi(\Delta n)$ is equivalent at multidegree $d$
to $\chi(n)^{\otimes C}$. In particular for $k=1$ the multiplier is $d!$.
\end{lemma}
\begin{proof}
On the positive Rees algebra $\mathcal R$ of Section~\ref{sec:scale}, set
\[
 T(Xt^a)=\left(X,a!\sum_{w(J)=a}X\epsilon_J\right)
                       \in\g\oplus\widetilde\g.           \tag{8}
\]
This is a rational Lie map. Indeed for each $U$ with $w(U)=a+b$ the
number of ordered disjoint decompositions $U=J\sqcup L$ with
$w(J)=a,w(L)=b$ is $\prod_i\binom{a_i+b_i}{a_i}$.
Multiplying by $a!b!$ gives $(a+b)!$. Thus
$[T(Xt^a),T(Yt^b)]=T([X,Y]t^{a+b})$ exactly. If $a+b$ is outside
the support, both sides vanish by $[\g_a,\g_b]\subseteq\g_{a+b}=0$.

Define $\mathfrak k_b=T(\bigoplus_{a\geq b,a\ne0}\g_a t^a)$,
$K_b=\exp\mathfrak k_b$. As before these give a descending, rational,
taut filtration on $K_0\subseteq G\times\widetilde G$, with lattice
$\Lambda=K_0\cap(\Gamma\times\widetilde\Gamma)$. They give the actual
comparison orbit
\[
 (g(n),\widetilde g(\Delta n))
       =\prod_a\exp\left(n^aT(X_at^a)\right).             \tag{9}
\]
No asserted equality of top weights is being used to generate a
non-descending filtration: every defining sum includes all $a\geq b$.
The top is exactly
\[
 K_d=\{(\exp X,\exp(CX\epsilon_{[D]})):X\in\g_d\}.       \tag{10}
\]
It is rational and central. The function on the comparison nilmanifold is
\[
 \mathcal F(x,y)=F(x)^{\otimes C}\otimes\overline{\widetilde F(y)}.
                                                        \tag{11}
\]
Its two vertical phases on (10) are $C\eta(\exp X)$ and
$-\widetilde\eta(\exp(CX\epsilon_{[D]}))=-C\eta(\exp X)$.
Therefore (11) descends to a Lipschitz function $F'$ on
$K_0/K_d$ with the projected lattice $\Lambda'$ and quotient filtration.
Writing $k_a'$ for the projected $\exp T(X_a t^a)$ gives the corrected
identity
\[
 \chi(n)^{\otimes C}\otimes\overline{\widetilde\chi(\Delta n)}
       =F'\left(\prod_a(k_a')^{n^a}\Lambda'\right).        \tag{12}
\]
The right side has top group zero, so it has support strictly below $d$.
For $k=1$, (11) contains $F^{\otimes d!}$ and (12) contains
$\chi^{\otimes d!}$. Dropping either factor would leave the uncancelled
vertical phase $(1-d!)\eta$ and would not prove the claimed equivalence.

The diagonal itself is a degree-$d$ nilcharacter: coarsen (5) by taking
the direct sum of components with $w(J)\geq a$ at index $a$. This is
descending, rational and taut, its top is still the $J=[D]$ component,
and (7) on the diagonal is polynomial for it. This verifies the exact
symbol domain used in (12), without assuming that arbitrary linear
maps preserve arbitrary multidegree filtrations.
\end{proof}

\section{Literal symmetry and division before construction}\label{sec:symmetry}
The map $\widetilde g$ is equivariant under block permutations:
$\widetilde g(\sigma z)=\sigma\widetilde g(z)$ after fixing the compatible
action convention. An arbitrary $\widetilde F$ need not be invariant,
so equivariance alone does not prove that $\widetilde\chi$ is symmetric.
We supply the missing transition explicitly.

Let $V$ be the finite-dimensional Hermitian target of $\widetilde F$ and
$m=|P|=C$. Regard $\Sym^m V$ as the symmetric subspace of $V^{\otimes m}$
with its induced norm, and write $\odot$ for symmetrised tensor product.
On the compact nilmanifold put
\[
 B(x)=\mathop{\bigodot}_{\sigma\in P}\widetilde F(\sigma x),
 \qquad F_{\rm sym}(x)=B(x)/\|B(x)\|.                    \tag{13}
\]
The product of nonzero vectors in the symmetric algebra is nonzero:
after a choice of basis this algebra is a polynomial ring over $\C$,
an integral domain. Hence $\|B(x)\|>0$ everywhere. Continuity and
compactness give a positive minimum, so (13) is Lipschitz. Reindexing
the finite product proves exact $P$-invariance. Since $P$ fixes the top
group, (13) has vertical frequency $m\widetilde\eta$.
Consequently
\[
 \chi_{\rm sym}(z)=F_{\rm sym}(\widetilde g(z)\widetilde\Gamma)
                                                        \tag{14}
\]
is a literally block-symmetric multidegree-$(1,\ldots,1)$ nilcharacter.
On a block diagonal, $\sigma\widetilde g(\Delta n)=\widetilde g(\Delta n)$.
All factors in (13) then equal the same unit vector $v$, and
$v^{\odot m}=v^{\otimes m}$ has norm one. Thus the diagonal of (14)
is exactly $\widetilde\chi(\Delta n)^{\otimes m}$ in the symmetric
subspace, not merely a proportional vector with an uncontrolled scalar.
Lemma~\ref{lem:raw} therefore gives
\[
 [\chi_{\rm sym}(\Delta n)]=Cm[\chi(n)].                 \tag{15}
\]

\begin{theorem}[E.9 and E.10 with a symmetric output]\label{thm:multi}
For a multidegree-$d$ nilcharacter $\chi$ there is a literally
block-symmetric multidegree-$(1,\ldots,1)$ nilcharacter $M(\chi)$
whose block diagonal is equivalent to $\chi$. The one-variable case
is the asserted symmetric E.9 output. The assertion holds on any
given limit subset after restriction.
\end{theorem}
\begin{proof}
If $d\ne0$, first use Corollary~\ref{cor:divide} with $Q=Cm=C^2$
to choose the specified representative $\rho$ with $[\chi]=Cm[\rho]$.
Apply (4)--(14) to this $\rho$, including its actual polynomial
coefficients, lattice, function and vertical frequency. Equation (15)
now identifies its output diagonal with $[\chi]$. We do not divide
the already symmetrised function and assume that its symmetry survives.
If $C=1$, all constructions and division remain valid. Coordinates
with $d_i=0$ have empty blocks; the normalised polynomial is independent
of them. If all $d_i=0$, the original sequence is a constant unit vector
and that vector on the zero-variable domain gives the assertion exactly.
\end{proof}

\section{The exact input to Section 13}\label{sec:join}
Let $s\geq2$ be fixed and let $\chi(h,n)$ be the multidegree-$(1,s-1)$
nilcharacter occurring immediately before published (13.4).
First choose $\rho$ with $[\chi]=s[\rho]$ by
Corollary~\ref{cor:divide}, in that multidegree symbol group.
Then set $\widetilde\chi=M(\rho)$ from Theorem~\ref{thm:multi}.
It has $s$ variables, multidegree $(1,\ldots,1)$ and literal symmetry
in its last $s-1$ variables. Its actual diagonal gives
\[
 [\chi(h,n)]_{\mathrm{Symb}_s({}^*\Z^2)}
   =s[\widetilde\chi(h,n,\ldots,n)]_{\mathrm{Symb}_s({}^*\Z^2)}.
                                                        \tag{16}
\]
Here is the comparison of symbol groups in (16). For support $a\leq(1,s-1)$,
give the same group the degree filtration generated by the $G_a$ with
$|a|\geq j$. Commutators satisfy the index law, rationality follows
from the rational generated-subgroup lemma, and the orbit is polynomial
by its monomial expansion. At total weight $s$ the only possible index
is $(1,s-1)$, so deleting that top index leaves degree at most $s-1$.
Thus both multidegree equivalences used above become degree-$s$
equivalences, with the same functions and the same input $(h,n)$.
This proves precisely (13.4), including its multiplier $s$.

The sequence $\Theta(n)=\widetilde\chi(n,\ldots,n)$ is a degree-at-most-$s$
nilcharacter by the same coarsening with the binary cube. Fixed linear
pullbacks $n\mapsto n+h$ preserve that degree filtration and preserve
lower-degree equivalence; one may either use the monomial expansion
or published Corollary B.5 and E.8(vi) in its degree-only form.
No full symmetry in all $s$ variables is asserted here. Section 13's
later symmetrisation argument is exactly what must address the first
variable. Its uses of multilinearity require the separately proved
replacement for Lemma 13.2; that result is not inferred from (16).
Likewise the subsequent bias and density arguments remain separate.

\section{Uniformity, limit interpretation, and audit boundary}\label{sec:limits}
All nilpotent groups, rational subspaces, finite tensor dimensions,
integer powers and permutation groups above are fixed by the original
geometric data, $d,k$, and the specified standard divisibility/scaling
integers. They do not depend on orbit coefficients or interval length.
The bounded-basepoint normalisation supplied by \cite{TAUT} has a uniform
Lipschitz constant. Subgroup restrictions, rational quotients and metrics
are fixed compact geometric maps. The lattice and the test function in
(6)--(7) are selected once from those fixed data. The positive minimum
in (13) is for that fixed test function on that fixed compact space,
not an infimum over a varying family of independently chosen functions.
For a varying original test function with bounded Lipschitz norm, the
raw comparison (11) varies, but the auxiliary $\widetilde F$ and (13)
can still be fixed because their frequency and group are fixed.

Every coefficient change is a finite BCH expression or a rational scalar
multiple in a fixed Lie algebra. Apply it coordinatewise to a representing
sequence to obtain a limit map. The finite formulas, exact identities,
positive fixed lower bound in (13), and uniform Lipschitz estimates pass
to the natural ultrapowers used in published Appendix A. Where geometric
data range over a standard finite list, perform this construction for
each member of that list. No selection depending on an unbounded interval
length or uncontrolled pointwise choice is introduced.
All vertical identities here have that natural-ultrapower meaning: they
hold for every limit point and every limit top-group element, by transfer
of the exact ordinary identity. An identity tested only on standard
points would not suffice; no such restriction is used here.

The affected-use audit in this unit covers published E.8(v),(vii),
E.9, E.10, and (13.4) through the definition of $\Theta$.
It also checks the literal symmetry used for the last $s-1$ variables
at those sites. It does not assert completeness of all corrections to
GTZ. The arbitrary-test-function symmetry transition and the factorial
omissions are preserved as findings with proposed repairs awaiting
independent mathematical review.

\begin{thebibliography}{9}
\bibitem{GTZ} B. Green, T. Tao and T. Ziegler,
An inverse theorem for the Gowers $U^{s+1}[N]$-norm,
Annals of Mathematics 176 (2012), 1231--1372.
DOI: 10.4007/annals.2012.176.2.11. Authoritative final PDF pinned in contracts.
\bibitem{TAUT} Project proof, Tautification and the shifted nilcharacter
construction, project author theta,
tautification-003, Theorem 1 and Corollary 2. Pending independent review.
\bibitem{LP13} Project proof, local Lemma 13.2, author
theta, lemma13-2-002, Section 2.
Pending independent review. Exact version pinned in contracts.
\end{thebibliography}
\end{document}
